\chapter{Le théorème de Pythagore}
\label{exercices:13}
\section*{Exercice 13.1 -- Une diagonale}
Un rectangle mesure $6$ sur $11\unite{cm}$. Calculer la diagonale exacte puis arrondie au millimètre.

\section*{Exercice 13.2 -- Une échelle}
Une échelle de $13\unite m$ touche un mur vertical et repose sur un sol horizontal à $5\unite m$ du mur. À quelle hauteur arrive-t-elle ? Quelle serait la hauteur si le pied se trouvait à $8\unite m$, avec la même longueur d'échelle ?

\section*{Exercice 13.3 -- Rectangle ou non ?}
Étudier les triangles de côtés $(6,8,10)$ puis $(6,8,11)$. Citer précisément le résultat utilisé dans chaque cas.

\section*{Exercice 13.4 -- Distance dans un repère}
Dans un repère orthonormé, $A(-2;1)$ et $B(4;9)$. Calculer $AB$ en construisant mentalement un triangle rectangle. Pourquoi le mot « orthonormé » est-il utile ?
